\section{Introduction}

Large language models (LLMs) excel at reasoning, coding, and multimodal tasks~\cite{olfactory2025,staticanalysis2025,egi2025}, but their hardware and deployment requirements often confine them to cloud infrastructure or high-end local systems, raising privacy, cloud dependence, and cost concerns~\cite{energytoken2025,qwqrouting2025,tigani2025cpu,aydin2025integrated}.

Small language models (SLMs) are attractive for local deployment and lower-cost, data-sovereign operation. Fine-tuning or distillation can also make them competitive on domain-specific tasks~\cite{xie2025fusing,zhu2024distilling,rethinkllm2025}. They pair naturally with routing, verification, and planning mechanisms~\cite{aydin2025integrated,qwqrouting2025,ji2025raprag}. However, SLMs still trail frontier LLMs on broad reasoning and tasks requiring extensive contextual knowledge~\cite{egi2025,dorfner2024comparing}, which has driven recent work on hybrid and agent-based architectures that combine small and large models for a better capability--efficiency trade-off~\cite{hao2024hybrid,shao2025dot,multiagent2025}.

This is primarily a systems-design question. Across domains, hybrid and multi-agent architectures aim to improve not only accuracy but also latency, cost, and computational efficiency through edge--cloud verification, expert routing, retrieval, and planner-based orchestration~\cite{hao2024hybrid,lee2024orchestrallm,alabbasi2024teleoracle,ji2025raprag,chatzistefanidis2025symbiotic,han2025chest,energytoken2025}. These developments motivate a closer examination of how collaborative architectures allocate work between small and large models, when monolithic LLM baselines remain necessary, and when orchestration provides the best balance between capability and efficiency.

Accordingly, this review examines when hybrid or multi-agent SLM--LLM architectures provide a better trade-off between capability and efficiency than monolithic LLM baselines, and under which domain conditions routing, retrieval, planning, and role specialization allow smaller models to contribute effectively.

Despite this growing interest, existing secondary studies do not address this question directly. The only PRISMA-guided predecessor reviews the SLM category broadly rather than collaborative SLM--LLM systems as a primary unit of analysis~\cite{slmreview2025prisma}; narrative surveys offer valuable conceptual breadth but neither treat monolithic LLMs as controlled baselines nor extract latency, cost, and energy as first-class outcome variables~\cite{wang2024slmsurvey,chen2025collabmech,wang2025collabperf}. The 2025--2026 wave of hybrid and multi-agent orchestration architectures is covered only partially by these works, leaving a gap in systematic, traceable evidence about when collaboration outperforms monolithic deployment under realistic engineering constraints.

The main contributions are fourfold: (1) a synthesis of 54 studies on hybrid SLM--LLM architectures, retrieval-centered collaboration, ensembles, and multi-agent workflows; (2) a review structure organized around four practical questions on performance, efficiency, orchestration, and domain suitability; (3) a three-pass qualitative coding scheme that links every synthesis claim to specific primary studies and a resulting analytical matrix mapping each extracted field to the research question(s) it addresses, distinguishing directly extracted data from analytical inference and thereby enabling independent verification of conclusions; and (4) a corpus-level characterisation showing that 72.2\% of the included work dates to 2025, that control-layer presence is the most consistent differentiator between high- and low-performing collaborative designs, and that energy remains systematically underreported relative to accuracy across the field.

The remainder of the paper is organized as follows. Section~\ref{sec:related} positions this review within prior secondary studies. Section~\ref{sec:methodology} details the search protocol, eligibility criteria, and qualitative coding procedure. Section~\ref{sec:findings} presents and synthesizes the results across the four research questions. Section~\ref{sec:threats} discusses threats to validity, and Section~\ref{sec:conclusion} concludes.
